\chapter{Dérivation - Partie 1}
\thispagestyle{fancy}

\begin{center}
	\includegraphics[scale=0.25]{2025_01_20_86f6bd542dfdbc66e53cg-1}
\end{center}

Le mot \guillemets{dérivé} vient du latin \guillemets{derivare} qui signifiait \guillemets{détourner un cours d'eau}. Le mot a été introduit par le mathématicien franco-italien Joseph Louis Lagrange (1736; 1813) pour signifier que cette nouvelle fonction dérive (au sens de "provenir") d'une autre fonction.

\section{Coefficient directeur d'une droite (Rappel)}
Le coefficient directeur de la droite (AB) est égal à :

$$
\frac{5-3}{4-1}=\frac{2}{3}
$$

Le coefficient directeur de la droite (CD) est égal à :

$$
\frac{-1-1}{6-2}=\frac{-2}{4}=-\frac{1}{2}
$$

\begin{center}
	\includegraphics[scale=0.25]{2025_01_20_86f6bd542dfdbc66e53cg-1(1)}
\end{center}

\section{Tangente à une courbe et nombre dérivé}

Soit A est un point appartenant à la courbe représentative d'une fonction $f$. On construit un réseau de sécantes à la courbe passant toutes par le point $A$ telle que le $2^{e}$ point d'intersection soit de plus en plus proche de $A$.\\
\includegraphics[scale=0.25]{2025_01_20_86f6bd542dfdbc66e53cg-1(2)}


On constate que la position limite des sécantes passant par le point $A$ est une droite \guillemets{touchant} la courbe au point A.\\
\includegraphics[scale=0.25]{2025_01_20_86f6bd542dfdbc66e53cg-2(1)}

\begin{definition}
	On appelle tangente à la courbe d'une fonction $f$ au point A , la droite limite d'un réseau de sécantes passant par $A$ et dont le $2^{e}$ point se rapproche de $A$.
\end{definition}

\begin{remarque}
	Géométriquement, la droite tangente à la courbe en A \guillemets{frôle} la courbe en A .
\end{remarque}

\section{Nombre dérivé}


\begin{definition}
	La tangente à la courbe au point A d'abscisse $a$ est la droite passant par A dont le coefficient directeur s'appelle le nombre dérivé de la fonction $f$ en $a$ et se note $f^{\prime}(a)$.\\
\includegraphics[scale=0.25]{2025_01_20_86f6bd542dfdbc66e53cg-2}
\end{definition}

\begin{propriete}
	L'équation réduite de la tangente à la courbe de la fonction $f$ en A est de la forme $y=f^{\prime}(a) x+p$ où $p$ est un nombre réel.
\end{propriete}

\begin{methodeTitre}{Déterminer graphiquement un nombre dérivé}
	\grasInMet{Enoncé :}
	On a représenté la fonction $f$ définie sur $\mathbb{R}$ par $f(x)=x^{2}+2 x-3$ et sa tangente au point d'abscisse 2.\\
	\includegraphics[scale=0.25]{2025_01_20_86f6bd542dfdbc66e53cg-3}\\
	\begin{enumMet}
		\item Déterminer graphiquement le coefficient directeur de la tangente à la courbe représentative de $f$ au point A et en déduire le nombre dérivé en 2 .\\
		\item Donner une équation de la tangente.\\
		\item En s'aidant de la calculatrice graphique, reproduire la courbe de la fonction $f$ et sa tangente au point A d'abscisse 2.
	\end{enumMet}
	
	\grasInMet{Solution :}
	
	\begin{enumMet}
		\item Le coefficient directeur de la tangente est égal à :
		
		$$
		\frac{6}{1}=6
		$$
		
		Le nombre dérivé de $f$ en 2 est égal à 6 .\\
		Ainsi, la tangente à la courbe représentative de $f$ au point A est la droite passant par A et de coefficient directeur 6.\\
		\includegraphics[scale=0.25]{2025_01_20_86f6bd542dfdbc66e53cg-4}\\
		\item Une équation de la tangente au pont A d'abscisse 2 est de la forme $y=6 x+p$.
		
		On a: $f(2)=2^{2}+2 \times 2-3=5$.\\
		Le point A a pour coordonnées $\binom{2}{5}$.\\
		Or, le point A appartient à la tangente donc ses coordonnées $\binom{2}{5}$ vérifient l'équation de la tangente : $y=6 x+p$.
		
		Donc $5=6 \times 2+p$\\
		Et donc $p=5-12=-7$\\
		Une équation de tangente à la courbe représentative de $f$ au point A d'abscisse 2 est : $y=6 x-7$.\\
		\item À l'aide de la calculatrice, il est possible de tracer la tangente à une courbe en un point.
	\end{enumMet}
\end{methodeTitre}





